\documentclass[a4paper]{article}     

\usepackage{amsfonts}
\usepackage{amsmath}
\usepackage{amssymb}
\usepackage{graphicx}
\usepackage{color}

\newcommand{\R}{\mathbb{R}}

\begin{document}

\noindent \large Auteur: Sadia Nazary \\
\noindent \large Studierichting: Informatica\\
\noindent \large Oefeningenreeks 7D nr. 1.3\\

\medskip

\normalsize

Gegeven:\\

$f(x,y)=e^{x^2-y}$\\

Gevraagd: het Taylorpolynoom van graad $3$ in $(0,0)$.\\

We berekenen alle partiele afgeleiden tot en met orde $3$ in $(0,0)$.\\

$f(0,0)=1$\\

Eerste orde:\\

$\dfrac{\partial f}{\partial x}=2xe^{x^2-y}$\\

$\dfrac{\partial f}{\partial x}(0,0)=0$\\

$\dfrac{\partial f}{\partial y}=-e^{x^2-y}$\\

$\dfrac{\partial f}{\partial y}(0,0)=-1$\\

Tweede orde:\\

$\dfrac{\partial^2 f}{\partial x^2}
=
(2+4x^2)e^{x^2-y}$\\

$\dfrac{\partial^2 f}{\partial x^2}(0,0)=2$\\

$\dfrac{\partial^2 f}{\partial x\partial y}
=
-2xe^{x^2-y}$\\

$\dfrac{\partial^2 f}{\partial x\partial y}(0,0)=0$\\

$\dfrac{\partial^2 f}{\partial y^2}
=
e^{x^2-y}$\\

$\dfrac{\partial^2 f}{\partial y^2}(0,0)=1$\\

Derde orde:\\

$\dfrac{\partial^3 f}{\partial x^3}
=
(12x+8x^3)e^{x^2-y}$\\

$\dfrac{\partial^3 f}{\partial x^3}(0,0)=0$\\

$\dfrac{\partial^3 f}{\partial x^2\partial y}
=
-(2+4x^2)e^{x^2-y}$\\

$\dfrac{\partial^3 f}{\partial x^2\partial y}(0,0)=-2$\\

$\dfrac{\partial^3 f}{\partial x\partial y^2}
=
2xe^{x^2-y}$\\

$\dfrac{\partial^3 f}{\partial x\partial y^2}(0,0)=0$\\

$\dfrac{\partial^3 f}{\partial y^3}
=
-e^{x^2-y}$\\

$\dfrac{\partial^3 f}{\partial y^3}(0,0)=-1$\\

Dus:\\

$T_3(x,y)
=
1-y+\dfrac{1}{2}(2x^2+y^2)
+
\dfrac{1}{6}(-6x^2y-y^3)$\\

$=
1-y+x^2+\dfrac{y^2}{2}-x^2y-\dfrac{y^3}{6}$\\

$\boxed{
T_3(x,y)
=
1-y+x^2+\dfrac{y^2}{2}-x^2y-\dfrac{y^3}{6}
}$\\

\begin{thebibliography}{999}
%\bibitem{a} Referentie
\end{thebibliography}

\end{document}
